\section{Composition with the exact unary Boolean indicator}
\label{sec:unary-transport}

The parent text \emph{The complete primewise coordinate system},
in its section \emph{Exact unary counts and simultaneous logical
constraints}, defines a rectangular unary code, the indicator
\(\chi_p\), and truth polynomials for Boolean operations.
We use those exact coordinates here.  The additional result is the
complete bijection to the bounded divisor boxes and the computation
of every simultaneous truth fibre on that one retained domain.
The classical torsor identities proved in that text are not needed
again for this composition.

Keep the full range, with these names distinguished from
the shear matrix \(J\) and from a general divisor bound:
\[
 J_{\rm code}=6h,\quad M_{\rm code}=9h,\quad
 \Lambda=2J_{\rm code}M_{\rm code},\quad
 N_{\rm code}=2J_{\rm code}M_{\rm code}^{\,2}.
\]
For this section only abbreviate \(J_{\rm code},M_{\rm code}\)
by \(J_0,M_0\).  The full rectangle consists of
\[
 0\le j<J_0,\quad 1\le A,B\le M_0,\quad
 \varepsilon\in\{0,1\},
\]
with shell \(a=3h+j+1,\ R=4j+3,\ S=pa\).
Its code and decoder are
\begin{align}
 c&=(A-1)\Lambda+2M_0j+2(B-1)+\varepsilon,
       \label{eq:unary-code}\\
 A&=1+\lfloor c/\Lambda\rfloor,\quad
 v=c-\Lambda\lfloor c/\Lambda\rfloor,\quad
 j=\lfloor v/(2M_0)\rfloor,\notag\\
 B&=1+\lfloor v/2\rfloor-M_0j,\qquad
 \varepsilon=v-2\lfloor v/2\rfloor. \notag
\end{align}
Successive division by \(\Lambda,2M_0,2\) proves these are
inverse for \(0\le c<N_{\rm code}\).

Write \([P]\) for the integer \(0\) or \(1\) according as the
specified predicate \(P\) is false or true.  Define
\[
 \kappa_p(c)=[AB\mid a]\,[\gcd(A,B)=1],\qquad
 \chi_p(c)=\kappa_p(c)\,[R\mid A+p^{1-\varepsilon}B].
\]
This is precisely the parent's indicator: \(4AB\mid p+4j+3\)
is equivalent to \(AB\mid a\), while its gcd-floor bits
equal these divisibility and coprimality bits.
Set \(\mathcal C_p^{\rm bud}=\{c:\kappa_p(c)=1\}\) and
\(\mathcal C_p^+=\{c:\chi_p(c)=1\}\).

For completeness the parent's integer truth-polynomial convention is
\[
 \neg b=1-b,\quad b\land d=bd,\quad b\lor d=b+d-bd,
\]
\[
 b\Rightarrow d=1-b+bd,\qquad
 b\Longleftrightarrow d=1-b-d+2bd.
\]
For \(b=0\), the last two expressions are \(1\) and
\(1-d\); for \(b=1\), they are both \(d\).
The first three follow by the same two substitutions.
Induction on the expression therefore gives its exact
zero-or-one truth value.  We use that convention for \(F\) below.

\begin{theorem}[Unary codes and one retained tagged valuation box]
\label{thm:unary-box}
There is an explicit bijection
\[
 \Theta:\mathcal B_p:=
 \coprod_{\substack{0\le j<6h\\\varepsilon\in\{0,1\}}}
 \{(j,\varepsilon)\}\times
 \prod_{\ell\mid a_j}\{0,\ldots,2v_\ell(a_j)\}
 \ \longrightarrow\ \mathcal C_p^{\rm bud}.
\]
For \((j,\varepsilon,\boldsymbol e)\), its forward map is
\[
 u=\prod_{\ell\mid a_j}\ell^{e_\ell},\quad
 g=\gcd(a_j,u),\quad A=a_j/g,\quad B=u/g,\quad D=g^2/u,
\]
followed by \eqref{eq:unary-code}.
Its inverse decodes \(j,A,B,\varepsilon\), sets
\[
 D=a_j/(AB),\qquad u=B^2D,\qquad
 \boldsymbol e=(v_\ell(u))_{\ell\mid a_j}.
\]
All fibres of these inverse maps are singletons.
The preimage \(\mathcal B_p^+=\Theta^{-1}(\mathcal C_p^+)\)
is exactly
\[
 R_j\mid 4u+p^\varepsilon,
\]
on the same box.  This is the exterior condition for \(\varepsilon=0\)
and the middle condition for \(\varepsilon=1\).
Every point in this preimage reconstructs the positive integers
\begin{equation}\label{eq:unary-integer-lift}
 C=\frac{A+p^{1-\varepsilon}B}{R_j},\qquad
 (a_j,y,z)=(ABD,p^\varepsilon ACD,pBCD).
\end{equation}
On its image, the inverse from the ordered triple retains its first
denominator, obtains the reduced pair
\[
 (m,n)=\operatorname{red}\big((R_jy-S_j)/S_j\big),
\]
and then recovers \(A=m,\ B=n/p^{1-\varepsilon}\), the tag
\(\varepsilon=0\) when \(p\mid n\) and \(\varepsilon=1\) otherwise,
and the displayed inverse code and exponents.
\end{theorem}
\begin{proof}
Unique prime factorization makes the exponent box bijective to
\(\Div(a_j^2)\), with the displayed valuation inverse.
Write \(a_j=gA,u=gB\) with \(\gcd(A,B)=1\).
Since \(u\mid a_j^2\), one has \(B\mid gA^2\), hence
\(B\mid g\).  Therefore \(D=g/B=g^2/u\) is a positive
integer and \(a_j=ABD,u=B^2D\).
In particular \(AB\mid a_j\) and \(A,B\le a_j\le9h\);
so the forward code belongs to the specified rectangle
and has \(\kappa_p=1\).
Conversely \(D=a_j/(AB)>0\) gives
\(u=B^2D\mid a_j^2\), with quotient \(A^2D\);
coprimality gives \(\gcd(a_j,u)=BD\).
It follows that the two maps recover \(A,B,D,u\) exactly.
Unique prime factorization and the code decoder recover every
remaining label, proving both inverse compositions.

The integers \(A,p,R_j\) satisfy
\(\gcd(A,R_j)=\gcd(p,R_j)=1\).
Multiplying \(4u+p^\varepsilon\) by \(A\) modulo \(R_j\),
and using \(4ABD=p+R_j\), gives
\[
 A(4B^2D+p^\varepsilon)
       \equiv pB+p^\varepsilon A
       =p^\varepsilon(A+p^{1-\varepsilon}B)\pmod{R_j}.
\]
Both multipliers are units, so this is an equivalence.
For \(\varepsilon=1\), \(4u+p\equiv4(u+a_j)\pmod{R_j}\)
and \(4\) is a unit, giving exactly the middle gate.

On the passing locus \(C,D>0\) are integers and
\[
 4ABCD=pC+A+p^{1-\varepsilon}B.
\]
With common denominator \(pABCD\), the three reciprocals
in \eqref{eq:unary-integer-lift} have these three
numerators.  Their sum is therefore \(4/p\).
Moreover direct substitution yields
\[
 R_jy-S_j=p^\varepsilon A^2D,\qquad
 \frac{R_jy-S_j}{S_j}=\frac{A}{p^{1-\varepsilon}B}.
\]
This fraction is reduced because \(\gcd(A,B)=1\) and
\(A\le a_j<p\).  Also \(B<p\), so the divisibility
of its denominator by \(p\) recovers the tag.
This proves the asserted inverse on the exact image,
including the retained ordering.  The exterior chart chooses
the orientation \((A,pB)\); its transpose is recovered by
exchanging \(y,z\), as in Section~1.
In the middle channel the source chart has pair \((A,B)\),
whereas the map \(\mu_a(u)\) below has pair \((B,A)\);
their exact correspondence is the coordinate exchange \(P\),
with the same exchange of the final two denominators.
No orientation is silently identified.
\end{proof}

\begin{theorem}[Simultaneous truth fibres, unions, and intersections]
\label{thm:unary-truth-fibres}
Let \(P_1,\ldots,P_t\) be any finite list of specified arithmetic
predicates decidable by exact integer operations on
\(\mathcal C_p^+\), including divisibility, congruences,
equalities and inequalities on the reconstructed coordinates.
Define the evaluation map on that finite domain by
\[
 \rho(c)=([P_1(c)],\ldots,[P_t(c)])\in\{0,1\}^t.
\]
For \(\sigma\in\{0,1\}^t\), the full fibre is
\[
 H_\sigma=\{c\in\mathcal C_p^+:\rho(c)=\sigma\},\qquad
 \widehat H_\sigma=
 \{\xi\in\mathcal B_p^+:\rho(\Theta(\xi))=\sigma\}.
\]
The restriction of \(\Theta\) is a bijection
\(\widehat H_\sigma\to H_\sigma\), with exactly the
integer lift and inverse of Theorem~\ref{thm:unary-box}.

For any finite Boolean expression \(\Phi\) in these predicates,
including any simultaneous list of biconditionals, put
\(T_\Phi=\{\sigma:\Phi(\sigma)=1\}\).
Its full solution set is the disjoint union
\[
 H_\Phi=\coprod_{\sigma\in T_\Phi}H_\sigma,\qquad
 \Theta^{-1}(H_\Phi)=
           \coprod_{\sigma\in T_\Phi}\widehat H_\sigma .
\]
In particular, for any \(r\ge1\) solution subsets \(X_1,\ldots,X_r\)
of this same coded domain,
\[
 \Theta^{-1}\Big(\bigcap_iX_i\Big)=\bigcap_i\Theta^{-1}(X_i),
 \qquad
 \Theta^{-1}\Big(\bigcup_iX_i\Big)=\bigcup_i\Theta^{-1}(X_i).
\]
The iterated fibre product
\(X_1\times_{\mathcal C_p^+}\cdots
\times_{\mathcal C_p^+}X_r\) over the inclusions is exactly
\(\{(c,\ldots,c):c\in\bigcap_iX_i\}\).
Its inverse diagonal projection has singleton fibres.
The union projection
\(\coprod_iX_i\to\bigcup_iX_i\), \((i,c)\mapsto c\),
has exact fibre \(\{(i,c):c\in X_i\}\).
A disjoint inverse representation of the union is obtained
by choosing the least such \(i\).

The computable exact cardinalities are
\begin{align}
 |H_\sigma|&=\sum_{c\in\mathcal C_p^+}
     \prod_{\sigma_i=1}[P_i(c)]
     \prod_{\sigma_i=0}(1-[P_i(c)]),\label{eq:truth-cardinality}\\
 |H_\Phi|&=\sum_{\sigma\in T_\Phi}|H_\sigma|
       =\sum_{\xi\in\mathcal B_p^+}
                       [\Phi(\rho(\Theta(\xi)))],\notag\\
 \left|\bigcup_{i=1}^rX_i\right|
     &=\sum_{\varnothing\ne I\subseteq\{1,\ldots,r\}}
           (-1)^{|I|+1}\left|\bigcap_{i\in I}X_i\right|.\notag
\end{align}
Equivalently, writing \(F(c)\) for the parent's truth
polynomial, the exact selected tagged count and weighted count are
\[
 \sum_{c=0}^{N_{\rm code}-1}\chi_p(c)F(c),\qquad
 W_F=\sum_{c=0}^{N_{\rm code}-1}
                  (2-\varepsilon(c))\chi_p(c)F(c).
\]
Predicates involving \(C,D,y,z\) are evaluated on the passing
domain where these coordinates exist.  They can be assigned any
fixed bits outside that domain, since the factor \(\chi_p=0\)
removes those codes.  No fractional candidate is treated as an
integer lift.  For \(F=1\), \(W_F=2Q_p\).
For a general selection invariant under middle complementation
\(u\mapsto a_j^2/u\), \(W_F/2\) counts the selected unordered
pairs; without that invariance it is only the displayed
half-weighted tagged count and can be nonintegral.
\end{theorem}
\begin{proof}
Each code has exactly one truth vector, so the fibres \(H_\sigma\)
are pairwise disjoint and partition \(\mathcal C_p^+\).
The inverse bijection \(\Theta^{-1}\) preserves membership
in every predicate evaluated at that same code; this proves
the fibre statement and both set identities.
A tuple in the fibre product has equal images under the
inclusions, which means all of its entries are literally
the same code.  Conversely any member of the intersection
gives that diagonal tuple.  The union fibre follows from
the definition of the coproduct, and the least-index rule
is defined precisely on the union.

The product in \eqref{eq:truth-cardinality} is one exactly
when every bit agrees with \(\sigma\); otherwise one of
its factors is zero.  Summing proves its count.  Disjointness
and the bijection give the other two expressions for
\(|H_\Phi|\).  At a point belonging to \(k>0\) of the
\(X_i\), the inclusion--exclusion sum contributes
\(\sum_{s=1}^k(-1)^{s+1}\binom{k}{s}=1\);
at a point in none it contributes zero.  The equality follows
by expanding \((1-1)^k\) and then summing over the finite domain.
The parent's Boolean truth polynomials evaluate to the same bit
\([\Phi]\) at each code, so their displayed counts agree with
these fibres.

Finally the exterior chart represents both ordered orientations
of one unordered pair and the middle charts represent the two
distinct complementary orientations.  The preceding proofs and
the fixed-point-free middle complementation of Section~1 give
\(W_1=2Q_p\).  An invariant middle selection retains both or
neither of those orientations; division by two then gives the
unordered count.  If one orientation alone is retained its
weight is one, proving both the limitation and the possible
nonintegrality asserted above.
\end{proof}

\begin{proposition}[Constructive CRT evaluation inside every common truth fibre]
\label{prop:unary-crt-fibre}
Every fibre of Theorem~\ref{thm:unary-truth-fibres} can be
computed with its complete integer lifts, using the bounded
divisibility CRT of Theorem~\ref{thm:divisor-crt} for its
positive affine/divisibility atoms and exact finite tests for
all its other atoms.  Noncoprime moduli, nonunit or zero
affine coefficients, negations, and biconditionals are retained.

More explicitly, in one fixed component \((j,\varepsilon)\),
write \(u=\prod_{\ell\mid a_j}\ell^{e_\ell}\).
Include the original bound \(u\mid a_j^2\) and the original
gate \(4u\equiv-p^\varepsilon\pmod{R_j}\).
For a truth vector \(\sigma\), collect every atom required true
of the forms
\[
 u\mid N_i,\qquad f_i\mid u,\qquad
 \alpha_i u\equiv\beta_i\pmod{L_i},
 \quad N_i,f_i,L_i>0,\quad\alpha_i,\beta_i\in\Z,
\]
where the displayed parameters are fixed in that component.
Let \(\mathcal L_\sigma\) be the complete finite list returned
by Theorem~\ref{thm:divisor-crt} on those conditions.
For each \(u\) on this list reconstruct its unique code by
Theorem~\ref{thm:unary-box} and retain it exactly when
\[
 [P_i(c)]=\sigma_i\quad(1\le i\le t).
\]
This gives precisely \(H_\sigma\) in that component, without
multiplicity.  Summing the retained list lengths over
\(j,\varepsilon,\sigma\in T_\Phi\) gives \(|H_\Phi|\).
Storing the codes, exponents, and \eqref{eq:unary-integer-lift}
gives all positive integer lifts and their singleton inverse
fibres, not only that cardinality.
\end{proposition}
\begin{proof}
Each actual member of \(H_\sigma\) satisfies its positive
collected atoms and the original bound and gate.  The full
biconditional in Theorem~\ref{thm:divisor-crt} therefore places
its unique \(u\) in \(\mathcal L_\sigma\).
It passes the final truth-vector test, proving that none is lost.
Conversely each retained list element satisfies all original
arithmetic conditions and every required true or false atom,
so its code is in exactly the stated fibre.
Unique factorization and Theorem~\ref{thm:unary-box} preclude
duplicates in a fixed component; the retained \(j,\varepsilon\)
labels and disjoint truth vectors preclude duplicates between
components or vectors.

For completeness this is a finite algorithm for arbitrary
specified decidable atoms, even when none has the collected
affine form: in that case the list is the bounded divisor
list satisfying the original gate, and every additional atom
is tested by its specified exact rule.  If an atom's coefficients
depend on \(u\) or on another decoded coordinate, that atom is
evaluated in this last step, rather than incorrectly treating
a variable coefficient as constant in CRT.
A false divisibility or congruence atom is likewise tested
on each same integer \(u\); it is not replaced by an independent
choice of a failing representative.  The preliminary gcd
compatibilities, full residue fibre \(r+L\Z\), monoid
coefficients, and retained valuation intervals are exactly
those proved in Theorem~\ref{thm:divisor-crt}.
These calculations produce the entire list or prove it empty;
no existence assumption is added.  This proves termination,
both inclusions, and the reconstruction and counting statements.
\end{proof}

The previously proved actual-shell obstruction \(p=37,a=18\)
makes the common-domain requirement concrete.  On
\(\Div(324)\), the two predicates \(5\mid4u+1\) and
\(7\mid4u+1\) have respectively the nonempty truth sets
\(\{1,6,36,81\}\) and \(\{12,54\}\).
Their intersection on that same divisor domain is empty.
The code bijection transports precisely these two finite
subsets, so it also transports the empty intersection.
The unrestricted integer CRT lift \(u=96\) remains outside
the valuation budget.  Neither taking a Cartesian product
of these local truth sets nor introducing separate
witnesses changes this computed fibre.
